\documentclass[10pt,a4paper]{amsart}
\usepackage[margin=19mm]{geometry}
\usepackage{amsmath,amssymb,mathtools}
\usepackage[scr=rsfs,cal=euler]{mathalfa}
\usepackage{xfrac}
\usepackage[unicode,hidelinks,bookmarks=false]{hyperref}
\newcommand{\ie}{i.e.\ }
\newcommand{\cf}{cf.\ }
\newcommand{\resp}{respectively\ }
\newcommand{\Subsection}{\subsection}
\providecommand{\nobreakdash}{\nobreak\hbox{-}}
\setlength{\emergencystretch}{2em}
\multlinegap0pt
\usepackage{amssymb}
\usepackage{xcolor}
\usepackage{tensor}
\usepackage[shortlabels]{enumitem}
\usepackage{mathtools}
\newtheorem{lem}{Lemma}
\newcommand{\alp}{\alpha}
\newcommand{\br}[1]{\overline{#1}}
\newcommand{\bt}{\beta}
\newcommand{\calA}{\mathcal A}
\newcommand{\calF}{\mathcal F}
\newcommand{\gmm}{\gamma}
\newcommand{\tht}{\theta}
\title{On the upper bound for the vorticity growth of bi-rotational Euler flows without swirl}
\author{Khakim Egamberganov, Deokwoo Lim}
\date{Source: arXiv:2607.27560v1 (CC BY 4.0)}
\begin{document}
\maketitle

\noindent\emph{Source and licence.} On the upper bound for the vorticity growth of bi-rotational Euler flows without swirl. Version arXiv:2607.27560v1 (CC BY 4.0), distributed by arXiv under the Creative Commons Attribution 4.0 International licence, http://creativecommons.org/licenses/by/4.0/, as recorded in the arXiv licence metadata for this exact version and shown on the abstract page https://arxiv.org/abs/2607.27560v1. Full licence notice: \texttt{licenses/arxiv-2607.27560-CC-BY-4.0.txt}.\par\smallskip

\noindent\textit{Editorial scope.} Counted unit: the article's lem lem:Frslow (source0.tex lines 1021--1027). The article's complete original proof of it is delivered with it (source0.tex lines 1029--1120, one \texttt{proof} environment of 92 lines). No mathematics is written, repaired or re-derived: the statement and the proof are the article's own lines, in the article's own notation, and no retained line is substituted or re-flowed.  Retained uncounted as the source-local context the proof needs: lines 918--929; lines 926--928.  Nothing was omitted from any retained span, so the package carries no omission record.  No retained line contains a citation, so the wrapper carries no bibliography and no bibliography entry.  No retained line contains a figure environment, an image-inclusion command or drawing code, so no image is delivered and the package declares no asset dependency.  Editorial formatting only, in addition: an AMS \texttt{amsart} preamble that restates the article's own declarations and macros used by the retained lines, namely the environments \texttt{lem} on separate counters, together with the standard packages those lines need.  The retained environments print in this document as Lemma 1 in order of appearance, whereas the article prints the same items with the numbering of its own full document.  The complete native source of the article as distributed in the arXiv e-print for this exact version is bundled unchanged as \texttt{sources/206353/source0.tex}.
	\begin{lem}\label{lem:calAdecay}
		Let $\alp\in \{
		0,1\}$, $\bt\geq0$, $\gmm>(\bt+1)/2$, and define
		\begin{equation}
			\calA_{\alp,\bt,\gmm}(\tau):=\int_{0}^{\pi}\frac{\cos^{\alp}\tht\sin^{\bt}\tht}{[\tau+2(1-\cos\tht)]^{\gmm}}d\tht,\quad %\calF_{\alp,\bt}(\tau):=\int_{0}^{\pi}\frac{\sin^{\bt}\tht}{[\tau+2(1-\cos\tht)]^{\alp}}d\tht,\quad 
			\tau>0.
		\end{equation}
		Then for any $\tau>0$, we have
		\begin{equation}\label{eq:calAdecay}
			0\leq\calA_{\alp,\bt,\gmm}(\tau)\lesssim_{\alp,\bt,\gmm} \min\lbrace \tau^{-(\gmm-(\bt+1)/2)}, \tau^{-(\gmm+\alp)}\rbrace.
		\end{equation}
	\end{lem}

		\begin{equation}
			0\leq\calA_{\alp,\bt,\gmm}(\tau)\lesssim_{\alp,\bt,\gmm} \min\lbrace \tau^{-(\gmm-(\bt+1)/2)}, \tau^{-(\gmm+\alp)}\rbrace.
		\end{equation}

	\begin{lem}\label{lem:Frslow}
		%Let . Then 
		For any $(r,s), (\br{r},\br{s})\in \Pi$, the kernels $F_{n,m}^{r}$ and $F_{n,m}^{s}$ satisfy
		\begin{equation}\label{eq:Frslow}
			|F_{n,m}^{r}(r,s,\br{r},\br{s})|, |F_{n,m}^{s}(r,s,\br{r},\br{s})|\lesssim_{n,m} \frac{1}{D}.%\min\left\{ \left(\frac{\br{r}}{r}\right)^{1/2}\frac{1}{D},\; \left(\frac{\br{s}}{s}\right)^{1/2}\frac{1}{D}\right\}.
		\end{equation} 
	\end{lem}

	\begin{proof}[Proof of Lemma \ref{lem:Frslow}]
		Throughout the proof, we often use the following inequality
		\begin{equation}\label{eq:basic}
			r, \br{r}\lesssim |r-\br{r}|+(r\br{r})^{1/2},\quad s, \br{s}\lesssim |s-\br{s}|+(s\br{s})^{1/2},
		\end{equation}
		which is obtained by
		\begin{equation}
			\begin{aligned}
				r&\leq r^{1/2}[|r-\br{r}|+\br{r}]^{1/2}\lesssim r^{1/2}|r-\br{r}|^{1/2}+(r\br{r})^{1/2}\leq \frac{r}{2}+\frac{|r-\br{r}|}{2}+(r\br{r})^{1/2},\\
				\br{r}&\leq \br{r}^{1/2}[|r-\br{r}|+r]^{1/2}\lesssim \frac{\br{r}}{2}+\frac{|r-\br{r}|}{2}+(r\br{r})^{1/2}.
			\end{aligned}
		\end{equation}
		First, we use the inequality \eqref{eq:basic} to get
		\begin{equation}
			\begin{aligned}
				&|F_{n,m}^{r}(r,s,\br{r},\br{s})|\lesssim \int_{0}^{\pi}\frac{\br{r}^{n}\br{s}^{m}\sin^{m-1}\br{\phi}_{1}|\br{s}-s\cos\br{\phi_{1}}|}{(r\br{r})^{d/2}}%\frac{\br{r}\br{s}|\br{s}-s\cos\br{\phi_{1}}|}{r^{2}\br{r}^{2}}
				\calA_{1,n-1,d/2}\left(\frac{D^{2}+2s\br{s}(1-\cos\br{\phi}_{1})}{r\br{r}}\right)d\br{\phi}_{1}\\
				\quad&=\int_{0}^{\pi}\frac{r\br{r}^{n+1}\br{s}^{m}\sin^{m-1}\br{\phi}_{1}|\br{s}-s\cos\br{\phi_{1}}|}{r^{d/2+1}\br{r}^{d/2+1}}\calA_{1,n-1,d/2}\left(\frac{D^{2}+2s\br{s}(1-\cos\br{\phi}_{1})}{r\br{r}}\right)d\br{\phi}_{1}\\
				\quad&\lesssim %\br{r}^{1/8}
				\int_{0}^{\pi}\frac{[|r-\br{r}|^{n+2}+(r\br{r})^{n/2+1}]\br{s}^{m}\sin^{m-1}\br{\phi}_{1}|\br{s}-s\cos\br{\phi_{1}}|}{r^{d/2+1}\br{r}^{d/2+1}}\calA_{1,n-1,d/2}\left(\frac{D^{2}+2s\br{s}(1-\cos\br{\phi}_{1})}{r\br{r}}\right)d\br{\phi}_{1}.
			\end{aligned}
		\end{equation}
		Then using the decay property
		\begin{equation}
			\calA_{1,n-1,d/2}(\tau)\lesssim_{n,m} \min\{ \tau^{-(m/2+1)},\tau^{-(d/2+1)}\}
		\end{equation}
		of the integral $ \calA_{1,n-1,d/2} $ from \eqref{eq:calAdecay} in Lemma \ref{lem:calAdecay}, we have
		\begin{equation}
			\begin{aligned}
				&|F_{n,m}^{r}(r,s,\br{r},\br{s})|\\
				\quad&\lesssim _{n,m}%\br{r}^{1/8}
				\int_{0}^{\pi}\frac{\br{s}^{m}\sin^{m-1}\br{\phi}_{1}|\br{s}-s\cos\br{\phi_{1}}|}{r^{d/2+1}\br{r}^{d/2+1}}\\
				&\qquad\cdot\left[|r-\br{r}|^{n+2}\left( \frac{r\br{r}}{D^{2}+2s\br{s}(1-\cos\br{\phi}_{1})}\right)^{d/2+1}+(r\br{r})^{n/2+1} \left( \frac{r\br{r}}{D^{2}+2s\br{s}(1-\cos\br{\phi}_{1})}\right)^{m/2+1}\right]d\br{\phi}_{1}\\
				\quad&\lesssim \int_{0}^{\pi}\frac{%\br{r}^{1/8}
					\br{s}^{m}\sin^{m-1}\br{\phi}_{1}}{[D^{2}+2s\br{s}(1-\cos\br{\phi}_{1})]^{(m+1)/2}}d\br{\phi}_{1}=\frac{\br{s}^{m}}{(s\br{s})^{(m+1)/2}}\calA_{0,m-1,(m+1)/2}\left(\frac{D^{2}}{s\br{s}}\right).
			\end{aligned}
		\end{equation}
		Then we use the inequality \eqref{eq:basic} and the decay
		\begin{equation}
			\calA_{0,m-1,(m+1)/2}(\tau)\lesssim_{m} \min\{ \tau^{-(m+1)/2},\tau^{-1/2}\}
		\end{equation}
		from Lemma \ref{lem:calAdecay} to obtain
		\begin{equation}
			\begin{aligned}
				&|F_{n,m}^{r}(r,s,\br{r},\br{s})|\lesssim \frac{[|s-\br{s}|^{m}+(s\br{s})^{m/2}]}{(s\br{s})^{(m+1)/2}}\calA_{0,m-1,(m+1)/2}\left(\frac{D^{2}}{s\br{s}}\right)%\int_{0}^{\pi}\frac{[|s-\br{s}|^{m}+(s\br{s})^{m/2}]\sin^{m-1}\br{\phi}_{1}}{[D^{2}+2s\br{s}(1-\cos\br{\phi}_{1})]^{(m+1)/2}}d\br{\phi}_{1}
				\\
				\quad&\lesssim_{m} \frac{%(\br{r}\br{s})^{1/8}
					1}{(s\br{s})^{(m+1)/2}}\left[|s-\br{s}|^{m}\left(\frac{s\br{s}}{D^{2}}\right)^{(m+1)/2}+(s\br{s})^{m/2}\left(\frac{s\br{s}}{D^{2}}\right)^{1/2}\right]\lesssim \frac{1%(\br{r}\br{s})^{1/8}
				}{D}.
			\end{aligned}
		\end{equation}
		Similarly, we use the inequality \eqref{eq:basic} to get
		\begin{equation}
			\begin{aligned}
				&|F_{n,m}^{s}(r,s,\br{r},\br{s})|\lesssim \int_{0}^{\pi}\frac{\br{r}^{n}\br{s}^{m}\sin^{n-1}\br{\tht}_{1}|\br{r}-r\cos\br{\tht_{1}}|}{(s\br{s})^{d/2}}%\frac{\br{r}\br{s}|\br{s}-s\cos\br{\phi_{1}}|}{r^{2}\br{r}^{2}}
				\calA_{1,m-1,d/2}\left(\frac{D^{2}+2r\br{r}(1-\cos\br{\tht}_{1})}{s\br{s}}\right)d\br{\tht}_{1}\\
				\quad&=\int_{0}^{\pi}\frac{\br{r}^{n}s\br{s}^{m+1}\sin^{n-1}\br{\tht}_{1}|\br{r}-r\cos\br{\tht_{1}}|}{s^{d/2+1}\br{s}^{d/2+1}}\calA_{1,m-1,d/2}\left(\frac{D^{2}+2r\br{r}(1-\cos\br{\tht}_{1})}{s\br{s}}\right)d\br{\tht}_{1}\\
				\quad&\lesssim %\br{r}^{1/8}
				\int_{0}^{\pi}\frac{\br{r}^{n}[|s-\br{s}|^{m+2}+(s\br{s})^{m/2+1}]\sin^{n-1}\br{\tht}_{1}|\br{r}-r\cos\br{\tht_{1}}|}{s^{d/2+1}\br{s}^{d/2+1}}\calA_{1,m-1,d/2}\left(\frac{D^{2}+2r\br{r}(1-\cos\br{\tht}_{1})}{s\br{s}}\right)d\br{\tht}_{1}.
			\end{aligned}
		\end{equation}
		Then we use the decay property
		\begin{equation}
			\calA_{1,m-1,d/2}(\tau)\lesssim_{n,m} \min\{ \tau^{-(n/2+1)},\tau^{-(d/2+1)}\},
		\end{equation}
		which gives us
		\begin{equation}
			\begin{aligned}
				&|F_{n,m}^{s}(r,s,\br{r},\br{s})|\\
				\quad&\lesssim %\br{r}^{1/8}
				\int_{0}^{\pi}\frac{\br{r}^{n}\sin^{n-1}\br{\tht}_{1}|\br{r}-r\cos\br{\tht_{1}}|}{s^{d/2+1}\br{s}^{d/2+1}}\\
				&\qquad\cdot\left[|s-\br{s}|^{m+2}\left( \frac{s\br{s}}{D^{2}+2r\br{r}(1-\cos\br{\tht}_{1})}\right)^{d/2+1}+(s\br{s})^{m/2+1} \left( \frac{s\br{s}}{D^{2}+2r\br{r}(1-\cos\br{\tht}_{1})}\right)^{n/2+1}\right]d\br{\tht}_{1}\\
				\quad&\lesssim \int_{0}^{\pi}\frac{%\br{r}^{1/8}
					\br{r}^{n}\sin^{n-1}\br{\tht}_{1}}{[D^{2}+2r\br{r}(1-\cos\br{\tht}_{1})]^{(n+1)/2}}d\br{\tht}_{1}=\frac{\br{r}^{n}}{(r\br{r})^{(n+1)/2}}\calA_{0,n-1,(n+1)/2}\left(\frac{D^{2}}{r\br{r}}\right).
			\end{aligned}
		\end{equation}
		Finally, we use the decay
		\begin{equation}
			\calA_{0,n-1,(n+1)/2}(\tau)\lesssim_{n} \min\{ \tau^{-(n+1)/2},\tau^{-1/2}\}
		\end{equation}
		to obtain
		\begin{equation}
			\begin{aligned}
				&|F_{n,m}^{s}(r,s,\br{r},\br{s})|\lesssim %(\br{r}\br{s})^{1/8}
				\frac{[|r-\br{r}|^{n}+(r\br{r})^{n/2}]}{(r\br{r})^{(n+1)/2}}\calA_{0,n-1,(n+1)/2}\left(\frac{D^{2}}{r\br{r}}\right)%\int_{0}^{\pi}\frac{[|s-\br{s}|^{m}+(s\br{s})^{m/2}]\sin^{m-1}\br{\phi}_{1}}{[D^{2}+2s\br{s}(1-\cos\br{\phi}_{1})]^{(m+1)/2}}d\br{\phi}_{1}
				\\
				\quad&\lesssim_{n} \frac{%(\br{r}\br{s})^{1/8}
					1}{(r\br{r})^{(n+1)/2}}\left[|r-\br{r}|^{n}\left(\frac{r\br{r}}{D^{2}}\right)^{(n+1)/2}+(r\br{r})^{n/2}\left(\frac{r\br{r}}{D^{2}}\right)^{1/2}\right]\lesssim \frac{1%(\br{r}\br{s})^{1/8}
				}{D}.
			\end{aligned}
		\end{equation}
	\end{proof}

\end{document}
